\chapter{Counting and Combinatorics}

There is a secret hiding in the first probability you ever computed. When you said
``a fair coin lands heads with probability $1/2$,'' or ``a die shows a four with
probability $1/6$,'' you were not really doing probability at all. You were
\emph{counting}. You counted the outcomes you cared about (one: heads), you counted
all the outcomes there were (two: heads, tails), and you divided. The whole edifice
of discrete probability rests on this one humble act.

That is the thesis of this chapter, and I want you to feel its weight before we start.
Whenever the outcomes of an experiment are \emph{equally likely} -- every card equally
probable to be drawn, every seating arrangement equally probable, every poker hand
equally probable -- the probability of an event is nothing more than
\[
  \P(\text{event}) \;=\; \frac{\text{number of outcomes in the event}}{\text{total number of outcomes}}.
\]
This is sometimes called the \emph{naive definition of probability}, and we will sharpen
it in the next chapter. But ``naive'' is a misleading name. The formula is exactly right
whenever the symmetry of the setup justifies it, and an astonishing fraction of real
problems -- cards, dice, lotteries, random samples, shuffles, hashing, password
strength -- have exactly that symmetry. The catch is the word \emph{number}. To use the
formula you must be able to count, and counting large structured sets is a genuine
mathematical art. That art is \emph{combinatorics}, and this chapter builds it up from
scratch.

\begin{focusbox}
By the end of this chapter you will be able to count: arrangements (permutations),
selections (combinations), things drawn with or without replacement, ordered or
unordered. You will know the binomial coefficient inside and out, build Pascal's
triangle and read its secrets, resolve the famous birthday paradox, wield
inclusion--exclusion to count overlapping things, and use ``stars and bars'' to count
the ways of distributing identical objects. Every technique will be tied back to the
counting-equals-probability idea above.
\end{focusbox}

\section{The Multiplication Principle}
\label{ch2:multiplication}

Everything in this chapter grows from a single seed, and the seed is so simple it feels
like a tautology. Suppose you are getting dressed and you own $3$ shirts and $4$ pairs of
trousers. How many distinct outfits can you assemble? For each of the $3$ shirts you may
pair any of the $4$ trousers, so the outfits come in $3$ groups of $4$, giving $3 \times 4 = 12$.

\begin{intuitionbox}
The picture to keep in your head is a \emph{tree}. From a single root you draw one branch
for each shirt; from the end of each shirt-branch you draw one branch for each pair of
trousers. The number of complete root-to-leaf paths is the number of outfits. Because
every shirt-branch sprouts the same number of trouser-branches, you do not add, you
\emph{multiply}.
\end{intuitionbox}

\begin{figure}[H]
\centering
\begin{tikzpicture}[level distance=14mm, sibling distance=10mm,
  every node/.style={font=\small}, grow'=right]
\node[circle,draw,fill=sourceblue!20] {start}
  child {node[circle,draw,fill=green!15] {$S_1$}
    child {node[fill=orange!12,draw,rounded corners] {$T_1$}}
    child {node[fill=orange!12,draw,rounded corners] {$T_2$}}
    child {node[fill=orange!12,draw,rounded corners] {$T_3$}}
    child {node[fill=orange!12,draw,rounded corners] {$T_4$}}
  }
  child {node[circle,draw,fill=green!15] {$S_2$}
    child {node[fill=orange!12,draw,rounded corners] {$T_1$}}
    child {node[fill=orange!12,draw,rounded corners] {$T_2$}}
    child {node[fill=orange!12,draw,rounded corners] {$T_3$}}
    child {node[fill=orange!12,draw,rounded corners] {$T_4$}}
  }
  child {node[circle,draw,fill=green!15] {$S_3$}
    child {node[fill=orange!12,draw,rounded corners] {$T_1$}}
    child {node[fill=orange!12,draw,rounded corners] {$T_2$}}
    child {node[fill=orange!12,draw,rounded corners] {$T_3$}}
    child {node[fill=orange!12,draw,rounded corners] {$T_4$}}
  }
;
\end{tikzpicture}
\caption{The multiplication principle as a tree. Three shirts, each branching into four
trousers, give $3 \times 4 = 12$ leaves, one per outfit.}
\label{ch2:fig-tree}
\end{figure}

\begin{theorem}[Multiplication Principle]
\label{ch2:thm-mult}
Suppose a procedure consists of $k$ successive stages, and suppose that
\begin{itemize}
  \item stage $1$ can be carried out in $n_1$ ways, and
  \item for each outcome of stage $1$, stage $2$ can be carried out in $n_2$ ways, and
  \item in general, for each outcome of the first $i-1$ stages, stage $i$ can be carried
        out in $n_i$ ways.
\end{itemize}
Then the whole procedure can be carried out in $n_1 \cdot n_2 \cdots n_k$ ways.
\end{theorem}

\begin{whybox}
Why is the product right, and not, say, the sum? Induct on the number of stages. With one
stage there are obviously $n_1$ ways. Now suppose the claim holds for $k-1$ stages, so the
first $k-1$ stages can be completed in $N = n_1 \cdots n_{k-1}$ ways. Fix any one of those
$N$ partial completions. The crucial hypothesis is that \emph{no matter which} of the $N$
we fixed, stage $k$ offers exactly $n_k$ continuations. So each of the $N$ partial paths
fans out into $n_k$ full paths, and these groups of $n_k$ never overlap (they start from
different partial paths). Total full paths: $N \cdot n_k = n_1 \cdots n_k$. The induction
closes.
\end{whybox}

The hypothesis I emphasized in the proof -- ``$n_i$ ways \emph{regardless of} the earlier
choices'' -- is the only thing that can trip you up. The \emph{number} of options at stage
$i$ must be constant; the options themselves may change.

\begin{example}[License plates]
\label{ch2:ex-plates}
A license plate has the form letter-letter-digit-digit-digit. How many plates are there
if repeats are allowed? Five stages: $26 \times 26 \times 10 \times 10 \times 10 =
676{,}000$. Now suppose no symbol may repeat. The letters: $26$ then $25$ (any letter but
the first). The digits: $10$, then $9$, then $8$. The count is $26 \times 25 \times 10
\times 9 \times 8 = 468{,}000$. Notice the option \emph{count} at the second letter stage
is $25$ no matter which first letter you chose -- so the multiplication principle still
applies cleanly even though \emph{which} $25$ letters are available changes.
\end{example}

\begin{example}[Subsets and the power of two]
\label{ch2:ex-subsets}
How many subsets does a set of $n$ elements have? Build a subset by walking through the
elements one at a time and, for each, deciding \emph{in} or \emph{out}. That is $n$ stages
of $2$ choices each, so $2^n$ subsets. This includes the empty set (everyone ``out'') and
the full set (everyone ``in''). A set of $10$ elements has $1024$ subsets; this exponential
blow-up is exactly why brute-force search over subsets becomes hopeless fast, and why we
will need cleverer counting.
\end{example}

\begin{happeningbox}
What is really happening is that \emph{independent choices compose by multiplication}.
This is the combinatorial shadow of a fact you will meet again and again in probability:
when events are independent, their probabilities multiply. The multiplication principle is
independence \emph{before} we have even mentioned chance. Keep that resonance in mind; in
the next chapter ``equally likely outcomes built from independent stages'' will turn
products of counts into products of probabilities.
\end{happeningbox}

\section{Permutations: Counting Arrangements}
\label{ch2:permutations}

An \emph{arrangement} is an ordering. In how many orders can $n$ distinct books be placed
on a shelf? Fill the leftmost slot first: $n$ choices. The next slot: $n-1$ books remain.
Then $n-2$, and so on down to $1$ for the last slot. By the multiplication principle the
number of orderings is
\[
  n \cdot (n-1) \cdot (n-2) \cdots 2 \cdot 1 \;=\; n!,
\]
read ``$n$ factorial.'' We define $0! = 1$, the empty arrangement; there is exactly one
way to arrange nothing, and as we will see this convention makes every later formula come
out right.

\begin{definition}[Factorial and $k$-permutations]
\label{ch2:def-perm}
For an integer $n \ge 0$, $n! = n(n-1)\cdots 1$ with $0! = 1$. The number of ways to choose
and arrange $k$ of $n$ distinct objects in order, written $P(n,k)$ or $(n)_k$, is
\[
  P(n,k) \;=\; n(n-1)\cdots(n-k+1) \;=\; \frac{n!}{(n-k)!}, \qquad 0 \le k \le n.
\]
\end{definition}

The product $n(n-1)\cdots(n-k+1)$ has exactly $k$ factors -- one per filled slot -- counting
down from $n$. The closed form $n!/(n-k)!$ is the same product with the unused tail
$(n-k)(n-k-1)\cdots 1$ canceling top and bottom. When $k = n$ the tail is $0! = 1$ and we
recover $P(n,n) = n!$, the full arrangement.

\begin{example}[Race podium]
\label{ch2:ex-podium}
Ten sprinters race; how many ways can gold, silver, bronze be awarded (no ties)?
Order matters -- gold then silver is different from silver then gold -- so this is
$P(10,3) = 10 \cdot 9 \cdot 8 = 720$.
\end{example}

\begin{trapbox}
The single most common error in counting is failing to ask: \textbf{does order matter?}
``Arrange three of ten runners on a podium'' (order matters: $P(10,3)=720$) is a different
question from ``choose three of ten runners for a relay pool'' (order does not matter, and
the answer will be $720/3! = 120$). Before you compute anything, decide whether two
selections that differ only in order should be counted once or twice. Get this wrong and
your count is off by a factor of $k!$ -- a mistake that scales catastrophically.
\end{trapbox}

\begin{historybox}
Factorials and arrangements are ancient. The Indian mathematician and astronomer
\emph{Bhaskara II} (12th century), and earlier Jain mathematicians, knew rules for
counting permutations and combinations; the medieval Hebrew scholar \emph{Abraham ibn Ezra}
computed combinations in an astrological context. The notation ``$n!$'' is far younger --
introduced by Christian Kramp in 1808 precisely because writing out the long product was
tedious. The \emph{idea} of counting arrangements long predates probability theory itself;
Cardano, Pascal, and Fermat inherited a rich counting tradition and pointed it at games of
chance.
\end{historybox}

\subsection{Permutations with repeated objects}

Suppose some objects are identical. How many distinct ``words'' can you make by rearranging
the letters of \texttt{MISSISSIPPI}? There are $11$ letters, so if all were distinct there
would be $11!$ arrangements. But the four \texttt{S}'s are interchangeable: any of their
$4!$ internal orderings produces the \emph{same} word, so we have overcounted by a factor
of $4!$. Likewise the four \texttt{I}'s overcount by $4!$ and the two \texttt{P}'s by $2!$.
The correct count is
\[
  \frac{11!}{4!\,4!\,2!\,1!} \;=\; 34650.
\]
This is the \emph{multinomial} count, and it embodies a principle worth naming.

\begin{intuitionbox}
\textbf{Overcount, then divide by the symmetry.} Few counting strategies are this general:
pretend everything is distinguishable and count the easy (larger) number, then divide by
the number of ways the indistinguishable items could be permuted among themselves without
changing anything. We will use this exact move to derive combinations in a moment.
\end{intuitionbox}

\section{Combinations: Counting Selections}
\label{ch2:combinations}

Now order stops mattering. A \emph{combination} is an unordered selection -- a subset. How
many $3$-person committees can be formed from $10$ people? If order mattered we computed
$P(10,3) = 720$ ordered selections. But a committee $\{A,B,C\}$ is the same committee no
matter the order we named its members, and there are $3! = 6$ orderings of any given trio.
So each committee was counted $6$ times among the $720$, and the true count is $720/6 = 120$.

\begin{definition}[Binomial coefficient]
\label{ch2:def-binom}
The number of $k$-element subsets of an $n$-element set is the \emph{binomial coefficient}
\[
  \binom{n}{k} \;=\; \frac{P(n,k)}{k!} \;=\; \frac{n!}{k!\,(n-k)!}, \qquad 0 \le k \le n,
\]
read ``$n$ choose $k$.'' By convention $\binom{n}{k} = 0$ when $k < 0$ or $k > n$.
\end{definition}

\begin{whybox}
The derivation is the overcount-then-divide move in its purest form. To build an ordered
selection of $k$ from $n$ you may first \emph{choose} the unordered set of $k$ ($\binom{n}{k}$
ways) and then \emph{order} it ($k!$ ways). By the multiplication principle the ordered count
factors as
\[
  P(n,k) \;=\; \binom{n}{k}\cdot k!.
\]
Solving for $\binom{n}{k}$ gives $P(n,k)/k! = n!/(k!(n-k)!)$. Every binomial coefficient is
therefore a count of ordered things with the order divided back out.
\end{whybox}

\begin{example}[Poker hands]
\label{ch2:ex-poker}
A standard deck has $52$ cards. A five-card hand is an unordered selection, so the number of
distinct hands is
\[
  \binom{52}{5} \;=\; \frac{52!}{5!\,47!} \;=\; \frac{52\cdot 51\cdot 50\cdot 49\cdot 48}
  {120} \;=\; 2{,}598{,}960.
\]
Now we can do real probability. How many hands are four-of-a-kind? Choose the rank that
appears four times ($13$ ways), take all four of its cards (one way), then choose the fifth
card from the remaining $48$ ($48$ ways): $13 \cdot 48 = 624$ hands. Because all
$\binom{52}{5}$ hands are equally likely,
\[
  \P(\text{four of a kind}) \;=\; \frac{624}{2{,}598{,}960} \;\approx\; 0.00024.
\]
Counting and probability are one motion: count the event, count the whole, divide.
\end{example}

\begin{figure}[H]
\centering
\begin{tikzpicture}
\begin{axis}[
  width=12cm, height=6.5cm,
  ybar, bar width=14pt,
  ymode=log, log origin=infty,
  ymin=1e-4, ymax=1,
  ylabel={probability (log scale)},
  symbolic x coords={4 of a kind, full house, flush, straight, 3 of a kind, two pair, one pair},
  xtick=data, x tick label style={rotate=35, anchor=east, font=\footnotesize},
  enlarge x limits=0.09,
  ymajorgrids, grid style={gray!18},
]
\addplot[fill=distblue!55, draw=distblue] coordinates {
  (4 of a kind, 0.00024) (full house, 0.001441) (flush, 0.001981)
  (straight, 0.003940) (3 of a kind, 0.021128) (two pair, 0.047539)
  (one pair, 0.422569)};
\end{axis}
\end{tikzpicture}
\caption{Probabilities of seven poker hands, each obtained by the same recipe: count the
hands of that type, divide by $\binom{52}{5}=2{,}598{,}960$. The vertical axis is
logarithmic because the rarities span more than three orders of magnitude -- four of a kind
is roughly $1760$ times rarer than one pair. (Here ``flush'' counts all five-same-suit hands
and ``straight'' all five-in-sequence hands; you will count two of these yourself in the
exercises.)}
\label{ch2:fig-poker}
\end{figure}

\begin{aestheticbox}
The notation $\binom{n}{k}$ is more telling than ``$nCk$'' because it reminds you the answer
is a single number attached to a pair $(n,k)$, and because it slots straight into the
binomial theorem
$(x+y)^n = \sum_{k=0}^n \binom{n}{k} x^k y^{n-k}$ that gave it its name. We will return to
that theorem; for now notice that the very same numbers count subsets \emph{and} expand
powers. When one quantity answers two unrelated-looking questions, mathematics is telling
you the questions were secretly the same.
\end{aestheticbox}

The same divide-by-the-whole recipe ranks every poker hand at once.
Figure~\ref{ch2:fig-poker} plots the probabilities of seven of them on a logarithmic
scale; the steep staircase down to four of a kind is a picture of why rare hands beat
common ones.

\subsection{Symmetry: choosing what to leave out}

A small identity carries a large idea:
\[
  \binom{n}{k} \;=\; \binom{n}{n-k}.
\]
Algebraically it is immediate from the formula, since swapping $k$ and $n-k$ leaves
$k!(n-k)!$ unchanged. But the \emph{reason} is prettier than the algebra. To choose $k$
objects to take is exactly to choose $n-k$ objects to leave behind; each selection of a
committee is simultaneously a selection of the non-committee. The two acts are the same act
seen from opposite sides, so they must have the same count. This is a \emph{bijective}
proof -- we exhibited a one-to-one correspondence rather than manipulating symbols -- and
bijective arguments are the soul of combinatorics. Whenever two counts are equal, ask
whether you can \emph{see} why by pairing the things up.

\section{The Binomial Coefficient and Pascal's Triangle}
\label{ch2:pascal}

Arrange the binomial coefficients in a triangle, row $n$ holding $\binom{n}{0}, \binom{n}{1},
\dots, \binom{n}{n}$:

\begin{figure}[H]
\centering
\begin{tikzpicture}[every node/.style={font=\small}]
\foreach \r [count=\ri from 0] in {
  {1},
  {1,1},
  {1,2,1},
  {1,3,3,1},
  {1,4,6,4,1},
  {1,5,10,10,5,1},
  {1,6,15,20,15,6,1},
}{
  \foreach \v [count=\ci from 0] in \r {
    \node at ({\ci*1.05 - \ri*0.525}, {-\ri*0.8})
      [draw, circle, inner sep=1.2pt, fill=sourceblue!10, minimum size=7mm] {\v};
  }
}
\draw[->,structurepurple,thick]
  ({2*1.05 - 5*0.525 - 0.45},{-5*0.8+0.35}) -- ++(0.7,0.55)
  node[midway,above,xshift=-6pt,structurepurple]{};
\node[structurepurple,font=\footnotesize] at (4.2,-3.0)
  {$\binom{5}{2}=\binom{4}{1}+\binom{4}{2}=4+6=10$};
\end{tikzpicture}
\caption{Pascal's triangle. Each entry is the sum of the two entries above it: the edges
are all $1$'s ($\binom{n}{0}=\binom{n}{n}=1$), and each interior number is born from the
pair diagonally above. Row $n$ sums to $2^n$.}
\label{ch2:fig-pascal}
\end{figure}

The defining miracle of the triangle is that each number is the sum of the two above it.
In symbols this is \emph{Pascal's rule}.

\begin{theorem}[Pascal's rule]
\label{ch2:thm-pascal}
For $1 \le k \le n$,
\[
  \binom{n}{k} \;=\; \binom{n-1}{k-1} + \binom{n-1}{k}.
\]
\end{theorem}

\begin{whybox}
You could prove this by grinding the factorials together over a common denominator, and it
works. But the \emph{combinatorial} proof tells you why it is true. We are counting
$k$-element subsets of $\{1,2,\dots,n\}$, which is $\binom{n}{k}$ by definition. Now sort
those subsets by a single question: \emph{does the subset contain the element $n$?}
\begin{itemize}
  \item If it \emph{does} contain $n$, the remaining $k-1$ members come from $\{1,\dots,n-1\}$:
        there are $\binom{n-1}{k-1}$ such subsets.
  \item If it \emph{does not} contain $n$, all $k$ members come from $\{1,\dots,n-1\}$:
        there are $\binom{n-1}{k}$ such subsets.
\end{itemize}
Every $k$-subset falls into exactly one case (it either has $n$ or it does not), so the two
counts add to the total. That is Pascal's rule, and the ``two entries above'' in the
picture are precisely these two cases.
\end{whybox}

This conditioning move -- split a count by whether some particular element is in or out -- is
a recurring theme, and it foreshadows \emph{conditioning} on an event, the engine of the
next chapters. Pascal's rule also gives a fast, multiplication-free way to compute the
coefficients: just add your way down the triangle.

\begin{historybox}
The triangle is named for Blaise Pascal, who in his 1654 \emph{Trait\'e du triangle
arithm\'etique} organized these numbers and used them to solve the ``problem of points'' (how
to fairly split the stakes of an interrupted game) in his famous correspondence with Pierre
de Fermat -- correspondence usually marked as the birth of probability theory. But the
triangle is much older: it appears in the work of the Persian mathematicians al-Karaji and
Omar Khayyam, in China as ``Yang Hui's triangle'' (13th century, citing earlier sources),
and in India centuries before. Pascal's contribution was not discovery but systematic use --
turning a numerical curiosity into a tool for reasoning about chance.
\end{historybox}

\subsection{Two identities you should be able to see}

\textbf{The row sum.} Adding across row $n$,
\[
  \sum_{k=0}^{n} \binom{n}{k} \;=\; 2^n.
\]
Why? The left side counts subsets of an $n$-set grouped by size; the right side counts all
subsets at once (in/out for each element, Example~\ref{ch2:ex-subsets}). Same set, two
counts, so they are equal. No algebra needed.

\textbf{The hockey-stick.} Summing down a diagonal,
\[
  \sum_{i=r}^{n} \binom{i}{r} \;=\; \binom{n+1}{r+1}.
\]
The picture (a straight diagonal of entries that ``bends'' at the bottom into the answer)
gives the rule its name. We will not belabor the proof, but it follows by repeatedly
applying Pascal's rule, and it is a handy closed form for sums that look intimidating.

\section{The Four Sampling Scenarios}
\label{ch2:four}

Almost every elementary counting problem is one of four flavors, indexed by two yes/no
questions. Are we sampling \emph{with replacement} (can an item be chosen more than once)?
Does \emph{order matter}? Here are all four, for choosing $k$ items from $n$.

\begin{figure}[H]
\centering
\renewcommand{\arraystretch}{1.8}
\begin{tabular}{l|c|c}
 & \textbf{Order matters} & \textbf{Order does not matter} \\
\hline
\textbf{Without replacement} & $\dfrac{n!}{(n-k)!}$ & $\dbinom{n}{k}$ \\
\hline
\textbf{With replacement} & $n^k$ & $\dbinom{n+k-1}{k}$ \\
\end{tabular}
\caption{The four ways to sample $k$ objects from $n$. The bottom-right cell -- ordered
\emph{no}, replacement \emph{yes} -- is the surprising one and is the subject of the
``stars and bars'' section below.}
\label{ch2:fig-four}
\end{figure}

Three of the four we have essentially seen. Ordered without replacement is the permutation
$P(n,k) = n!/(n-k)!$ (each pick removes an option, and order is recorded). Unordered without
replacement is the combination $\binom{n}{k}$. Ordered with replacement is the simplest of
all: $k$ independent stages of $n$ choices each, so $n^k$ -- this is the license-plate count,
the password count, the ``rolls of a die'' count. The fourth, unordered \emph{with}
replacement, is genuinely subtle and we devote a full section to it.

\begin{example}[Three flavors of ``three from a deck'']
\label{ch2:ex-three}
\emph{Deal $3$ cards, recording order, no card reused:} $52 \cdot 51 \cdot 50 = 132600$.
\emph{Deal $3$ cards as an unordered hand:} $\binom{52}{3} = 22100 = 132600/3!$.
\emph{Draw a card, note it, replace, reshuffle, repeat $3$ times, recording order:}
$52^3 = 140608$. The same deck, three different questions, three different counts -- and
each is the denominator of a different probability calculation.
\end{example}

\section{The Birthday Problem}
\label{ch2:birthday}

Here is a question that has been startling people for a century. In a room of just $23$
people, the probability that \emph{some two of them share a birthday} is already better
than even. Most people guess you would need around $180$ people (half of $365$). The gap
between intuition and truth is so large that the result is called the \emph{birthday
paradox} -- though, as with all good paradoxes, nothing is actually contradictory; our
intuition is simply miscalibrated.

\begin{paradoxbox}
\textbf{The birthday paradox.} Assume $365$ equally likely birthdays and ignore leap years
and twins. With $23$ people in a room, the chance that at least two share a birthday is
about $50.7\%$. With $50$ people it is $97\%$; with $70$ people it is $99.9\%$. How can so
few people suffice? The resolution: you are not comparing your birthday to one other
person's. You are comparing \emph{every pair} of people, and the number of pairs grows
quadratically. With $23$ people there are $\binom{23}{2} = 253$ pairs -- already comfortably
more than half of $365$ -- and each pair is a fresh chance for a collision.
\end{paradoxbox}

Now compute it. The trick is to count the \emph{complement}. Asking for ``at least
one shared birthday'' invites messy overlaps (maybe two pairs match, maybe three people all
match). The opposite event -- ``all $n$ birthdays distinct'' -- is clean: it is just an
ordered sampling without replacement. The first person can have any of $365$ birthdays, the
second any of the remaining $364$, and so on, while the total number of birthday assignments
(with replacement, since people may repeat) is $365^n$. So
\[
  \P(\text{all distinct})
  \;=\; \frac{365 \cdot 364 \cdots (365 - n + 1)}{365^n}
  \;=\; \prod_{i=0}^{n-1}\Bigl(1 - \frac{i}{365}\Bigr),
\]
and the probability of at least one match is $1$ minus that.

\begin{walkbox}
Before reading on, predict the shape of the curve. Each factor $(1 - i/365)$ is slightly
less than $1$, and you are multiplying $n$ of them. With many near-$1$ factors multiplied
together, the product drops faster than you would think -- a fact captured by
$\prod (1 - i/365) \approx \exp\!\bigl(-\sum i/365\bigr) = \exp\!\bigl(-\tfrac{n(n-1)}{2 \cdot 365}\bigr)$.
The exponent has $n(n-1)/2 = \binom{n}{2}$ in it -- the number of pairs! That is the
paradox in one line: the collision probability is governed by the \emph{number of pairs},
which is quadratic in $n$.
\end{walkbox}

\begin{figure}[H]
\centering
\begin{tikzpicture}
\begin{axis}[
  width=12cm, height=7cm,
  xlabel={number of people $n$}, ylabel={$\P(\text{at least one shared birthday})$},
  xmin=0, xmax=80, ymin=0, ymax=1.02,
  grid=both, grid style={gray!18},
  legend pos=south east, legend style={font=\small},
  samples=81, domain=1:80,
]
% closed-form approximation 1 - exp(-n(n-1)/730), accurate to <1% in this range
\addplot[densityred, thick, smooth] {1 - exp(-x*(x-1)/730)};
\addlegendentry{$\P(\ge 1\text{ match})$}
\addplot[dashed, gray] coordinates {(0,0.5) (80,0.5)};
\addplot[dashed, structurepurple] coordinates {(23,0) (23,1.02)};
\node[structurepurple, font=\footnotesize, anchor=west] at (axis cs:24,0.27)
  {$n=23$};
\node[font=\footnotesize, anchor=west] at (axis cs:2,0.55) {$50\%$ line};
\end{axis}
\end{tikzpicture}
\caption{The birthday curve (closed-form approximation $1 - e^{-n(n-1)/730}$). It crosses
$50\%$ near $n=23$ and is essentially certain by $n \approx 60$. The steep rise is the
signature of the quadratic number of pairs.}
\label{ch2:fig-birthday}
\end{figure}

\begin{example}[The numbers]
\label{ch2:ex-birthday-numbers}
For $n=23$, the exact product gives $\P(\text{all distinct}) = 0.4927$, so
$\P(\text{match}) = 0.5073$ -- just over half. For $n=30$ it is $0.706$; for $n=50$ it is
$0.970$; for $n=57$ it has passed $0.99$. The crossover at $23$ is the headline, but the
\emph{speed} of the climb afterwards is just as striking.
\end{example}

\begin{applicationbox}
\textbf{The birthday attack in cryptography.} This is not a parlor trick. A hash function
maps data to one of $N$ possible fingerprints; a ``collision'' is two different inputs with
the same fingerprint, and collisions can break digital signatures. The birthday math says
you expect a collision after hashing only about $\sqrt{N}$ random inputs, not $N$ -- because,
again, it is the number of \emph{pairs} that matters. To make collisions infeasible a hash
therefore needs roughly \emph{twice} as many output bits as you might naively demand: a
$128$-bit security target needs a $256$-bit hash. The phrase ``birthday bound'' is standard
in cryptography, and it traces directly to the count $\binom{n}{2}$ we did above. The same
$\sqrt{N}$ scaling appears in hash-table load analysis and in distributed systems that
assign random IDs.
\end{applicationbox}

\section{Inclusion--Exclusion}
\label{ch2:incexc}

The multiplication and combination tools count things that fit cleanly into stages or
subsets. But many counts involve \emph{overlap}: ``how many integers up to $100$ are
divisible by $2$ or $3$?'' If you add the multiples of $2$ ($50$ of them) to the multiples
of $3$ ($33$), you have double-counted everything divisible by both $2$ and $3$, i.e. by
$6$ ($16$ of them). Subtract those once and the count is $50 + 33 - 16 = 67$.

\begin{intuitionbox}
Picture two overlapping circles. To find the area of the union you add the two circle areas
and then \emph{subtract the lens} where they overlap, because you counted it twice. With
three circles you add the three, subtract the three pairwise overlaps, then add back the
triple overlap (which you removed one time too many). The alternating add--subtract--add
pattern is inclusion--exclusion.
\end{intuitionbox}

\begin{theorem}[Inclusion--Exclusion]
\label{ch2:thm-incexc}
For finite sets $A_1, \dots, A_n$,
\[
  \abs[\big]{A_1 \cup \cdots \cup A_n}
  \;=\; \sum_{i} \abs{A_i}
  \;-\; \sum_{i<j} \abs{A_i \cap A_j}
  \;+\; \sum_{i<j<k} \abs{A_i \cap A_j \cap A_k}
  \;-\; \cdots
  \;+\; (-1)^{n+1}\abs[\big]{A_1 \cap \cdots \cap A_n}.
\]
The same identity holds with $\abs{\cdot}$ replaced by $\P(\cdot)$, giving the
inclusion--exclusion formula for the probability of a union.
\end{theorem}

\begin{whybox}
The cleanest proof tracks a single element $x$ and asks: \emph{how many times does $x$ get
counted on the right-hand side?} Suppose $x$ lies in exactly $m \ge 1$ of the sets. Then $x$
is counted $\binom{m}{1}$ times in the singles, $\binom{m}{2}$ times in the pairs (once for
each pair of the $m$ sets containing it), $\binom{m}{3}$ in the triples, and so on. Its net
contribution to the right side is
\[
  \binom{m}{1} - \binom{m}{2} + \binom{m}{3} - \cdots
  \;=\; -\sum_{j=1}^{m} (-1)^{j}\binom{m}{j}
  \;=\; 1 - \sum_{j=0}^{m} (-1)^{j}\binom{m}{j}
  \;=\; 1 - 0 \;=\; 1,
\]
where we used the alternating row-sum $\sum_{j=0}^{m}(-1)^j\binom{m}{j} = (1-1)^m = 0$ for
$m \ge 1$ (the binomial theorem at $x=-1$, $y=1$). So every element of the union is counted
exactly once on the right, which means the right side equals $\abs{A_1\cup\cdots\cup A_n}$.
Elements in \emph{no} set contribute $0$ on both sides. Done.
\end{whybox}

\begin{example}[The hat-check problem / derangements]
\label{ch2:ex-derange}
$n$ guests check identical-looking hats; the absent-minded attendant returns them at
random. What is the probability that \emph{no one} gets their own hat back? Let $A_i$ be the
event that guest $i$ \emph{does} get their own hat. We want $\P(\text{none}) = 1 -
\P(A_1 \cup \cdots \cup A_n)$. By symmetry $\P(A_i) = 1/n$ (guest $i$'s hat is equally likely
to land anywhere), $\P(A_i \cap A_j) = \frac{(n-2)!}{n!} = \frac{1}{n(n-1)}$, and in general
the intersection of any $k$ of them has probability $(n-k)!/n!$. There are $\binom{n}{k}$
such $k$-fold intersections, and $\binom{n}{k}\frac{(n-k)!}{n!} = \frac{1}{k!}$. Inclusion--exclusion gives
\[
  \P(A_1\cup\cdots\cup A_n) = \sum_{k=1}^{n} (-1)^{k+1}\frac{1}{k!},
  \qquad\text{so}\qquad
  \P(\text{no match}) = \sum_{k=0}^{n}\frac{(-1)^k}{k!}.
\]
\end{example}

\begin{paradoxbox}
\textbf{The hat-check limit.} As $n \to \infty$ the no-match probability is
$\sum_{k\ge 0}(-1)^k/k! = e^{-1} \approx 0.3679$. The startling part: the answer barely
depends on $n$. Whether $5$ guests or $5$ million guests check their hats, the probability
that the random return gives \emph{nobody} their own hat hovers near $37\%$. Adding more
guests does not wash out the chance of a total mismatch, because each new guest brings both
a new chance to match \emph{and} a new hat to be mismatched, and these effects almost
perfectly cancel. The constant $e$ -- the hero of calculus -- shows up unbidden in a pure
counting problem.
\end{paradoxbox}

\begin{connectionbox}
That appearance of $e^{-1}$ is your first sighting of a bridge we will cross repeatedly:
the link between combinatorics and \emph{analysis}. The series $\sum (-1)^k/k!$ is the
Taylor expansion of $e^x$ evaluated at $x=-1$ -- the very series the companion calculus
book derives. Throughout probability, finite counting identities have limits that are
calculus objects: factorials smooth into the Gamma function, binomial coefficients smooth
into the Beta function, and the binomial distribution we build from $\binom{n}{k}$ will, in
the limit, become the bell curve via the Central Limit Theorem. Counting is discrete
probability; calculus is its continuous shadow.
\end{connectionbox}

\section{Stars and Bars}
\label{ch2:starsbars}

We saved the strangest cell of Figure~\ref{ch2:fig-four} for last: counting
\emph{unordered selections with replacement}. Equivalently -- and this is the framing that
unlocks it -- how many ways can we distribute $k$ identical objects into $n$ distinct boxes?

Concretely: you buy $k = 7$ identical doughnuts from a shop with $n = 3$ varieties. A
purchase is determined entirely by \emph{how many of each variety} you take, say $(3,2,2)$
or $(7,0,0)$ -- order of buying is irrelevant and you may repeat a variety. How many
purchases are possible?

\begin{intuitionbox}
Represent a purchase as a row of $7$ stars (the doughnuts) separated into $3$ groups by
$2$ bars (the dividers between varieties). For instance
\[
  \star\star\star \mid \star\star \mid \star\star
  \qquad\text{means } (3,2,2),
\]
\[
  \star\star\star\star\star\star\star \mid\, \mid
  \qquad\text{means } (7,0,0).
\]
Every purchase is one arrangement of $7$ stars and $2$ bars in a row; every such
arrangement is one purchase. To count the purchases, just count the arrangements of stars
and bars!
\end{intuitionbox}

\begin{figure}[H]
\centering
\begin{tikzpicture}[font=\large]
\foreach \x/\sym/\col in {
  0/$\star$/sourceblue, 0.55/$\star$/sourceblue, 1.1/$\star$/sourceblue,
  1.75/$\mid$/trapred,
  2.4/$\star$/sourceblue, 2.95/$\star$/sourceblue,
  3.6/$\mid$/trapred,
  4.25/$\star$/sourceblue, 4.8/$\star$/sourceblue}{
  \node[\col] at (\x,0) {\sym};
}
\draw[decorate,decoration={brace,amplitude=4pt},gray] (-0.2,-0.35) -- (1.3,-0.35)
  node[midway,below,font=\footnotesize,black]{$3$};
\draw[decorate,decoration={brace,amplitude=4pt},gray] (2.2,-0.35) -- (3.15,-0.35)
  node[midway,below,font=\footnotesize,black]{$2$};
\draw[decorate,decoration={brace,amplitude=4pt},gray] (4.05,-0.35) -- (5.0,-0.35)
  node[midway,below,font=\footnotesize,black]{$2$};
\end{tikzpicture}
\caption{Stars and bars for the purchase $(3,2,2)$: seven stars (doughnuts) and two bars
(dividers) in a row of nine positions. Each distinct arrangement of stars and bars is one
purchase, so the count is the number of ways to place the $2$ bars among $9$ slots,
$\binom{9}{2}$.}
\label{ch2:fig-starsbars}
\end{figure}

A row has $7 + 2 = 9$ positions, and we choose which $2$ of them hold bars (the rest are
stars). That is $\binom{9}{2} = 36$ (Figure~\ref{ch2:fig-starsbars}). In general:

\begin{theorem}[Stars and bars]
\label{ch2:thm-stars}
The number of ways to write a nonnegative integer $k$ as an ordered sum of $n$ nonnegative
parts $x_1 + x_2 + \cdots + x_n = k$ -- equivalently, the number of ways to place $k$
identical objects into $n$ distinct boxes, equivalently the number of unordered samples of
size $k$ from $n$ types with replacement -- is
\[
  \binom{n+k-1}{k} \;=\; \binom{n+k-1}{n-1}.
\]
\end{theorem}

\begin{whybox}
A configuration is a row of $k$ stars and $n-1$ bars: the $n-1$ bars cut the row into $n$
(possibly empty) groups, and the group sizes are exactly the $x_i$. There are $k + (n-1)$
total symbols, and a configuration is fixed once we decide which $n-1$ positions are bars --
that is $\binom{n+k-1}{n-1}$ -- or equivalently which $k$ positions are stars, which is
$\binom{n+k-1}{k}$ (the two are equal by the choosing-what-to-leave-out symmetry). The map
from configurations to solutions $(x_1,\dots,x_n)$ is a bijection: each star-and-bar row
yields exactly one tuple, and each tuple yields exactly one row. Bijection means equal
counts.
\end{whybox}

\begin{trapbox}
Stars and bars counts \emph{nonnegative} solutions (boxes may be empty). If every box must
be nonempty -- every variety bought at least once, every $x_i \ge 1$ -- first hand one
object to each box ($n$ objects spent), then distribute the remaining $k - n$ freely:
$\binom{(k-n)+n-1}{n-1} = \binom{k-1}{n-1}$. Forgetting which regime you are in (boxes
allowed empty vs. forced nonempty) is a classic slip. Always ask: can a part be zero?
\end{trapbox}

\begin{example}[Distributing identical gifts]
\label{ch2:ex-gifts}
In how many ways can $10$ identical candies be split among $4$ children (some children may
get none)? Stars and bars with $k=10$, $n=4$: $\binom{10+4-1}{4-1} = \binom{13}{3} = 286$.
If every child must get at least one candy, give one to each first and distribute the
remaining $6$: $\binom{6+4-1}{4-1} = \binom{9}{3} = 84$.
\end{example}

\begin{happeningbox}
Why does ``unordered with replacement'' equal stars and bars? An unordered sample with
replacement does not care which doughnut you grabbed first; it cares only about the
\emph{multiplicities} -- how many of each type. A list of multiplicities $(x_1,\dots,x_n)$
that sums to $k$ is exactly a stars-and-bars configuration. So the seemingly exotic
bottom-right cell of the sampling table is the same object as the down-to-earth ``how many
of each'' question. This is also why this count is $\binom{n+k-1}{k}$ and \emph{not}
$\binom{n}{k}$: with replacement the pool effectively grows, and the algebra records that
with the $+k-1$ in the top.
\end{happeningbox}

\section{Why Counting Is the Backbone of Discrete Probability}
\label{ch2:backbone}

Here, finally, is the payoff promised at the start. Whenever a finite experiment has equally
likely outcomes, every probability question becomes two counting questions: count the
favorable outcomes, count all outcomes, divide. The techniques of this chapter are
therefore not a warm-up \emph{before} probability; for equally-likely models they \emph{are}
probability.

\begin{formalbox}
\textbf{The naive model, stated.} Let the sample space $\Omega$ be a finite set of outcomes,
all assumed equally likely. For any event $A \subseteq \Omega$,
\[
  \P(A) \;=\; \frac{\abs{A}}{\abs{\Omega}}.
\]
Every fact you need follows. Nonnegativity: $\abs{A} \ge 0$. Normalization:
$\P(\Omega) = \abs{\Omega}/\abs{\Omega} = 1$. Additivity for disjoint events:
$\abs{A \cup B} = \abs{A} + \abs{B}$ when $A \cap B = \varnothing$, so the probabilities
add. These are exactly the axioms of probability we develop in the next chapter -- here
they are theorems about counting.
\end{formalbox}

\begin{example}[Two dice, the long way and the right way]
\label{ch2:ex-twodice}
Roll two fair dice. What is $\P(\text{sum} = 7)$? The sample space is the $6 \times 6 = 36$
ordered pairs (multiplication principle), all equally likely. The sum-$7$ outcomes are
$(1,6),(2,5),(3,4),(4,3),(5,2),(6,1)$ -- six of them. So $\P = 6/36 = 1/6$. The reason a
naive count over \emph{unordered} pairs would mislead you here is that unordered pairs are
\emph{not} equally likely (you can roll a $\{3,4\}$ two ways but a $\{3,3\}$ only one way).
The art is choosing a sample space whose outcomes truly are equally likely, and ordered
pairs are that space. \textbf{Choose the equally-likely space first; then count.}
\end{example}

\begin{aestheticbox}
Counting is the backbone for a deeper reason than convenience. The hardest part of applied
probability is rarely the arithmetic -- it is building an accurate model of the situation:
What are the outcomes? Are they equally likely? Does order matter? Can things repeat? Those
four questions are the very ones combinatorics forces you to answer precisely. Master
counting and you have mastered the discipline of stating a probability problem cleanly,
which is most of the battle. Everything later -- random variables, expectation,
distributions -- is built on sample spaces, and sample spaces are counted things.
\end{aestheticbox}

\begin{labbox}
Pick a small experiment of your own -- drawing $2$ cards, flipping $4$ coins, rolling $3$
dice -- and do the following before computing any probability. (1) Write down the sample
space explicitly enough to count it. (2) Argue why its outcomes are equally likely (or fix
the space until they are). (3) Express your event as a subset and count it with one of this
chapter's tools. (4) Only then divide. Doing this slowly a few times builds the reflex that
makes hard problems easy.
\end{labbox}

\section{Exercises}
\label{ch2:exercises}

\begin{exercise}[Warm-up counts]
\label{ch2:exr-warmup}
(a) How many $4$-digit PINs are there if digits may repeat? If they may not?
(b) How many ways can $7$ people sit in a row? Around a round table (rotations counted as
the same)?
(c) How many distinct rearrangements are there of the letters in \texttt{BOOKKEEPER}?
\end{exercise}

\begin{exercise}[Order or not]
\label{ch2:exr-order}
A class has $20$ students. (a) In how many ways can a President, Vice-President, and
Treasurer be chosen (distinct roles)? (b) In how many ways can a committee of $3$
(no roles) be chosen? (c) Explain the exact numerical relationship between your two answers
and why it is what it is.
\end{exercise}

\begin{exercise}[Poker probabilities]
\label{ch2:exr-poker}
Using the $\binom{52}{5} = 2{,}598{,}960$ hands, count and then find the probability of:
(a) a flush (all five cards the same suit, including straight flushes);
(b) exactly one pair (one rank twice, the other three cards of three distinct other ranks);
(c) compare your one-pair probability to the four-of-a-kind probability from
Example~\ref{ch2:ex-poker} and comment.
\end{exercise}

\begin{exercise}[Birthday variations]
\label{ch2:exr-birthday}
(a) Using the approximation $\P(\ge 1\text{ match}) \approx 1 - e^{-n(n-1)/730}$, estimate
the smallest $n$ giving probability at least $0.5$, and compare with the exact answer $23$.
(b) Suppose instead birthdays are spread over a planet with $D$ equally likely ``days.''
Show that the number of people needed for an even-chance collision grows like
$\sqrt{D}$, and state the constant in front of $\sqrt{D}$.
\end{exercise}

\begin{exercise}[Inclusion--exclusion in the wild]
\label{ch2:exr-incexc}
How many integers from $1$ to $1000$ are divisible by at least one of $2$, $3$, or $5$?
Then find how many are divisible by \emph{none} of them. (Hint: count the divisible ones
with inclusion--exclusion and subtract from $1000$.)
\end{exercise}

\begin{exercise}[Stars and bars]
\label{ch2:exr-stars}
(a) How many nonnegative integer solutions does $x_1 + x_2 + x_3 + x_4 = 12$ have?
(b) How many with every $x_i \ge 1$? (c) How many with $x_1 \ge 2$ and the rest $\ge 0$?
(Hint for (c): pre-assign the required minimum, then count freely.)
\end{exercise}

\begin{exercise}[A bijective proof]
\label{ch2:exr-bijection}
Prove the identity $\sum_{k=0}^{n} \binom{n}{k} = 2^n$ \emph{combinatorially} -- that is,
by describing a single set that both sides count -- rather than by algebra. Then prove
$\binom{n}{k} = \binom{n}{n-k}$ the same way, exhibiting an explicit bijection between the
two families of subsets.
\end{exercise}

\begin{exercise}[Vandermonde's identity]
\label{ch2:exr-vandermonde}
Prove Vandermonde's identity
$\binom{m+n}{r} = \sum_{j=0}^{r}\binom{m}{j}\binom{n}{r-j}$
by a counting argument. (Hint: you have $m$ red balls and $n$ blue balls and want to choose
$r$ in total; condition on how many of the chosen are red.) Then deduce the special case
$\sum_{k=0}^{n}\binom{n}{k}^2 = \binom{2n}{n}$.
\end{exercise}

\begin{exercise}[Derangement check]
\label{ch2:exr-derange}
Using Example~\ref{ch2:ex-derange}, compute the exact probability that no guest gets their
own hat for $n = 3, 4, 5$ guests, and verify that the values are already very close to
$e^{-1} \approx 0.3679$. How large must $n$ be for the answer to be within $0.001$ of
$e^{-1}$?
\end{exercise}

\begin{exercise}[Modeling care]
\label{ch2:exr-model}
A bag has $3$ red and $2$ blue marbles. You draw $2$ without replacement. A student argues:
``There are three unordered outcomes -- two red, two blue, one of each -- so each has
probability $1/3$.'' Explain precisely why this reasoning is wrong, identify the correct
equally-likely sample space, and compute the right probability of ``one of each.''
\end{exercise}
